\begin{ouverture}[courbe]{7}{Proportionnalité, rapports et ratios}{Partie 2 · Algèbre et modèles}
{Quatrième proportionnelle · Égalité de rapports · Ratios · Partages}
\objectif{Reconnaître une proportionnalité et calculer une quatrième proportionnelle.}
\objectif{Comparer deux nombres ou deux grandeurs par un rapport ou un ratio.}
\objectif{Résoudre des problèmes de partage proportionnel.}
\prerequis{Proportionnalité, tableaux et graphiques (Cinquième), fractions.}
\end{ouverture}
\pagestyle{otMain}\essentiel\fontsize{11}{15}\selectfont
\phantomsection\label{cours-7-1}
\essSection{1}{Coefficients et égalité des rapports}
\begin{definition}Deux grandeurs sont proportionnelles si l’on passe des valeurs de la première à celles de la seconde en multipliant toujours par un même nombre, le coefficient de proportionnalité. Quand les deux grandeurs sont de natures différentes, ce coefficient a une unité : 2,50 €/kg, 80 km/h.\end{definition}
\begin{propriete}[Égalité des rapports]Si $(a\,;b)$ et $(c\,;d)$ sont deux couples de valeurs correspondantes de grandeurs proportionnelles, avec $b\ne0$ et $d\ne0$, alors $\frac ab=\frac cd$. Pour $b\ne0$ et $d\ne0$, l’égalité $\frac ab=\frac cd$ équivaut à $a\times d=b\times c$ : les produits en croix sont égaux.\end{propriete}
\begin{exemple}4 cahiers identiques coûtent 10 €. Prix de 7 cahiers ? Retour à l’unité : un cahier coûte 2,50 €, donc 7 cahiers coûtent $7\times2{,}50=17{,}50$ €. Produits en croix : $\frac4{10}=\frac7p$ donne $4p=70$, donc $p=17{,}50$ €.\end{exemple}
\begin{attention}Un tarif de 5 € fixes plus 2 € par article n’est pas proportionnel : doubler le nombre d’articles ne double pas le prix. Une proportionnalité se représente par des points alignés avec l’origine. Le produit en croix n’est qu’une méthode : linéarité et retour à l’unité restent valables.\end{attention}
\clearpage
\phantomsection\label{cours-7-2}
\essSection{2}{Rapports, ratios et partages}
\begin{definition}Le rapport de deux nombres, ou de deux grandeurs de même nature, $a$ et $b$ ($b\ne0$) est le quotient $\frac ab$ : il permet de les comparer. On emploie aussi le mot ratio. Les nombres $a$ et $b$ sont dans le ratio $3:5$ lorsque $\frac a3=\frac b5$ ; $a$, $b$, $c$ sont dans le ratio $2:3:7$, ou proportionnels à 2, 3, 7, lorsque $\frac a2=\frac b3=\frac c7$.\end{definition}
\begin{savoirfaire}[Partager selon un ratio]Pour partager 56 dans le ratio $3:5$, on compte $3+5=8$ parts égales. Une part vaut $56:8=7$. Les deux quantités sont $3\times7=21$ et $5\times7=35$. Contrôle : $21+35=56$ et $\frac{21}3=\frac{35}5=7$.\end{savoirfaire}
\label{fig-cours-7-2-cours-ratio}\figureGeo{cours-ratio}

\begin{exemple}[Une échelle est un rapport]À l’échelle $1:500$, aussi notée $\frac1{500}$, 1 cm sur le plan représente 500 cm en réalité : on compare deux longueurs exprimées dans la même unité.\end{exemple}
\begin{attention}L’ordre compte : le ratio $3:5$ n’est pas le ratio $5:3$. Dans un partage selon $3:5$, la première part représente $\frac38$ du total, et non $\frac35$.\end{attention}

\clearpage\rubrique{Découverte}
\begin{blocLivre}{ex-7-1}\begin{exercice}[1]{Passer à l’unité}
Six bouteilles identiques coûtent 9 euros. Combien coûtent dix bouteilles, sans frais fixes ?
\end{exercice}\end{blocLivre}
\begin{blocLivre}{ex-7-2}\begin{exercice}[1]{Compléter un tableau}
Le prix est proportionnel à la masse. Compléter le tableau.\begin{center}\renewcommand{\arraystretch}{1.35}\begin{tabular}{c|ccc}\toprule
Masse (kg) & 2 & 5 & 8\\\midrule
Prix (euros) & 7 & ? & ?\\
\bottomrule\end{tabular}\end{center}
\end{exercice}\end{blocLivre}
\begin{blocLivre}{ex-7-3}\begin{exercice}[1]{Des produits en croix}
Les couples $(6\,;15)$ et $(8\,;x)$ sont des valeurs correspondantes de deux grandeurs proportionnelles. Calculer $x$ avec l’égalité des produits en croix, puis vérifier avec le coefficient.
\end{exercice}\end{blocLivre}
\begin{blocLivre}{ex-7-4}\begin{exercice}[1]{Un ratio dans une classe}
Dans une classe de 28 élèves, il y a 12 filles. Donner le rapport du nombre de filles au nombre de garçons, puis le ratio filles : garçons sous forme simplifiée.
\end{exercice}\end{blocLivre}
\begin{blocLivre}{ex-7-5}\begin{exercice}[1]{Lire une échelle}
Sur un plan à l’échelle $1:500$, un mur mesure 3,6 cm. Quelle longueur réelle représente-t-il, en mètres ?
\end{exercice}\end{blocLivre}

\rubrique{Entraînement}
\begin{blocLivre}{ex-7-6}\begin{exercice}[2]{Une quatrième proportionnelle}
Huit pièces identiques pèsent 14,4 kg. Quelle est la masse de cinq pièces ? Donner deux méthodes.
\end{exercice}\end{blocLivre}
\begin{blocLivre}{ex-7-7}\begin{exercice}[2]{Un tarif non proportionnel}
Un service facture 4 euros fixes puis 3 euros par heure. Calculer pour 1 h et 2 h. Le prix est-il proportionnel au temps ?
\end{exercice}\end{blocLivre}
\begin{blocLivre}{ex-7-8}\begin{exercice}[2]{Un graphique}
Le graphique représente un prix en fonction d’une masse. Reconnaître la proportionnalité et lire le prix de 3 kg.
\label{fig-ex-7-8}\figureGeo{proportion-prix}
\end{exercice}\end{blocLivre}
\begin{blocLivre}{ex-7-9}\begin{exercice}[2]{Partager selon un ratio}
Ana et Ben se partagent 72 euros dans le ratio $4:5$. Représenter le partage par un schéma en barres, puis calculer les deux parts.
\label{fig-ex-7-9}\figureGeo{ratio-72}
\end{exercice}\end{blocLivre}
\begin{blocLivre}{ex-7-10}\begin{exercice}[2]{Trois nombres proportionnels}
Les nombres $a$, $b$, $c$ sont proportionnels à 2, 3 et 7, et $a+b+c=60$. Calculer $a$, $b$ et $c$.
\end{exercice}\end{blocLivre}

\rubrique{Problèmes}
\begin{blocLivre}{ex-7-11}\begin{exercice}[3]{Une recette et des lots}
Une recette pour 6 personnes demande 450 g de farine. On la prépare pour 14 personnes. La farine est vendue en paquets de 500 g. Combien de paquets faut-il ?
\end{exercice}\end{blocLivre}
\begin{blocLivre}{ex-7-12}\begin{exercice}[3]{Une part connue}
Trois amis se partagent des cartes dans le ratio $4:3:2$. Le deuxième en reçoit 18. Combien les deux autres en reçoivent-ils ? Combien y a-t-il de cartes en tout ?
\end{exercice}\end{blocLivre}
\begin{blocLivre}{ex-7-13}\begin{exercice}[3]{Un arc de cercle}
Un cercle de rayon 6 cm porte un arc de mesure $75^\circ$. Calculer la longueur de cet arc, exactement puis au dixième de cm.
\label{fig-ex-7-13}\figureGeo{arc-75}
\end{exercice}\end{blocLivre}
\begin{blocLivre}{ex-7-14}\begin{exercice}[3]{Un mélange de peinture}
Une teinte s’obtient en mélangeant peinture bleue et peinture blanche dans le ratio $2:5$. On dispose de 3 L de bleu. Quelle quantité de blanc faut-il ? Quel volume de mélange obtient-on ? Quelle fraction du mélange est bleue ?
\end{exercice}\end{blocLivre}
\begin{blocLivre}{ex-7-15}\begin{exercice}[3]{Comparer deux recettes}
La recette A utilise 250 g de farine pour 0,5 L de lait ; la recette B 400 g de farine pour 0,75 L de lait. Laquelle contient le plus de farine par litre de lait ? Avec B, quelle masse de farine faut-il pour 1,2 L de lait ?
\end{exercice}\end{blocLivre}

\clearpage
